\def\stat{borisov-bosov}





\def\tit{СТРАТЕГИЯ ИССЛЕДОВАНИЙ И~РАЗРАБОТОК\\ В~ОБЛАСТИ ИСКУССТВЕННОГО ИНТЕЛЛЕКТА 


IV:\\ ГОСУДАРСТВЕННАЯ ПОЛИТИКА КНР}





\def\titkol{Стратегия исследований и~разработок в~области искусственного интеллекта 


IV} %:~Государственная политика КНР}





\def\aut{А.\,В.~Борисов$^1$, А.\,В.~Босов$^2$, Д.\,В.~Жуков$^4$}





\def\autkol{А.\,В.~Борисов, А.\,В.~Босов, Д.\,В.~Жуков}





\titel{\tit}{\aut}{\autkol}{\titkol}





\index{Борисов А.\,В.}


\index{Босов А.\,В.}


\index{Жуков Д.\,В.}


\index{Borisov A.\,V.}


\index{Bosov A.\,V.}


\index{Zhukov D.\,V.}





%{\renewcommand{\thefootnote}{\fnsymbol{footnote}}


%\footnotetext[1]


%{Публикация выполнена при поддержке Программы стратегического академического лидерства РУДН и~при 


%финансовой поддержке РФФИ (проекты 19-07-00933 и~20-07-01064).}} 





\renewcommand{\thefootnote}{\arabic{footnote}}


\footnotetext[1]{Федеральный исследовательский центр <<Информатика и~управление>> Российской академии наук, 


\mbox{ABorisov@ipiran.ru}}


\footnotetext[2]{Федеральный исследовательский центр <<Информатика и~управление>> Российской академии наук, 


\mbox{AVBosov@ipiran.ru}}


\footnotetext[3]{Федеральный исследовательский центр <<Информатика и~управление>> 


Российской академии наук, \mbox{DZhukov@ipiran.ru}}





 \label{st\stat}





 \thispagestyle{headings}


 


 \vspace*{-12pt} 








  





        \Abst{Статья завершает цикл работ, посвященных анализу влияния государственного 


управления на эффективность проведения исследований и~разработок в~об\-ласти искусственного 


интеллекта (AI R\&D). Четвертая часть цик\-ла представляет результат анализа влияния на область 


AI R\&D государственной политики Китайской народной республики (КНР). Представлены 


краткие сведения о~на\-уч\-но-тех\-ни\-че\-ском потенциале КНР и~стратегических документах 


КНР в~части развития AI R\&D, принятых с~2017~г. Описана структура государственного 


и~частного финансирования AI R\&D в~КНР. Дано краткое описание организации AI R\&D в~об\-ласти 


обороны и~безопасности. Представлены заключительные замечания, в~том чис\-ле различия 


в~подходах к~стратегическому государственному планированию в~об\-ласти AI R\&D, принятых 


в~различных государствах.}


        


        \KW{искусственный интеллект; сис\-те\-ма распознавания лиц; система  


наблю\-де\-ния\!/\!сле\-же\-ния; умный го\-род/безопас\-ный город}





\DOI{10.14357\!/\!08696527220102} 





\vskip 10pt plus 9pt minus 6pt





 \thispagestyle{headings}


 


 \vspace*{-20pt}


        


\section{Введение}


     


     Цикл работ~[1--3], который завершает статья, посвящен анализу 


государственного влияния на сферу исследований и~разработок в~об\-ласти 


искусственного интеллекта. Предыдущие час\-ти были посвящены 


хронологии развития данной об\-ласти в~\mbox{СССР} и~Российской Федерации 


и~сравнительному анализу наукометрических показателей публикационной 


активности по AI R\&D за последние 20~лет в~России и~мире~\cite{1-b, 2-b}, 


а также исследованию стратегических документов AI R\&D  


в~США~\cite{3-b}. Эти сведения видятся <<начальными условиями>> 


реализации российской <<Национальной стратегии развития искусственного 


интеллекта на период до 2030~года>>~\cite{4-b}. Окончательное дополнение 


их со\-сто\-ит в~рас\-смот\-ре\-нии стратегических планов научных исследований, 


разработок и~внед\-ре\-ния технологий искусственного интеллекта (AI) в~Китайской 


народной республике. Внимание к~КНР в~контексте представленного 


цикла работ объясняется сле\-ду\-ющи\-ми общеизвестными фактами:


     \begin{itemize}


\item КНР делит лидирующие позиции с~США в~сфере высоких технологий 


и~даже обгоняет их в~некоторых областях;


\item КНР, став региональным лидером, стремится к~мировому 


экономическому, политическому и~военному первенству;


\item политическая и~экономическая система КНР, выбравшей свой 


собственный <<китайский путь развития>>, существенно отличается как от 


социалистической, так и~от капиталистической сис\-те\-мы и~в~некоторых 


областях демонстрирует преимущества;


\item восточный менталитет накладывает свой отпечаток на стратегическое 


мыш\-ле\-ние, и~соответствующие подходы должны изучаться.


\end{itemize}





     Работа организована следующим образом. Раздел~2 содержит краткие 


сведения о~на\-уч\-но-тех\-ни\-че\-ском потенциале КНР и~стратегических 


государственных документах в~части развития AI R\&D, принятых с~2017~г. 


Раздел~3 посвящен структуре государственного и~частного финансирования 


AI R\&D в~КНР. Организация AI R\&D в~об\-ласти обороны и~безопас\-ности 


КНР изложена в~разд.~4. Заключительные замечания ко всему 


представленному циклу работ приведены в~разд.~5.


     





\section{Государственный подход к~AI R\&D в~КНР}


     


     Согласно американским материалам, КНР занимает первое место по 


публикациям в~научных журналах~\cite{5-b} по тематике искусственного 


интеллекта. Летом 2017~г.\ Госсовет КНР принял документ <<План 


разработки искусственного интеллекта нового поколения>>~[6, 7] (далее~--- 


План). Точные траты на AI в~КНР не указаны, однако некоторые <<реперные 


точки>> Китая позволяют оценить финансовые и~людские затраты в~этой 


об\-ласти. На 2017~г.\ в~КНР было более 80~млн ученых, т.\,е.\ более 5,8\% 


населения. Сто семьдесят миллионов человек, более 12\% населения, имеют высшее 


образование. В~2016~г.\ объем продаж на\-уч\-но-тех\-ни\-че\-ской 


продукции превысил 1~трлн юаней, это около 140~млрд долларов США. В~этом 


же году на на\-уч\-но-ис\-сле\-до\-ва\-тель\-ские работы (НИР) 


и~на\-уч\-но-ис\-сле\-до\-ва\-тель\-ские и~опыт\-но-кон\-струк\-тор\-ские 


работы (НИОКР) в~КНР было затрачено около~2,1\% валового внут\-рен\-не\-го продукта (ВВП)~--- 


1,544~млрд юаней, это более 216~млн долларов, из которых~78\% со\-ста\-ви\-ли 


част\-ные инвестиции.


     


     В~качестве основных принципов Плана указаны: достижение лидерства 


в~технологиях; системный подход; ориентация на рынок; открытость. 


В~документе определены стратегические цели и~этапы их достижения.


     


     \textbf{Первый этап ({до 2020~г.})}: обеспечение соответствия общего 


уровня технологий и~применения AI в~КНР передовому общемировому 


уровню:


     \begin{itemize}


\item разработка теории и~технологий обработки больших массивов 


разнородных данных, группового и~гибридного интеллекта, автономных 


интеллектуальных систем;


\item выведение индустрии AI КНР на первые общемировые позиции, 


создание ключевых предприятий по разработке AI, расширение финансовой 


базы индустрии AI с~150~млрд до 1~трлн юаней;


\item оптимизация условий НИР\!/\!НИОКР в~об\-ласти AI, насыщение отрасли 


высокопрофессиональными кадрами, разработка комплекса законодательных 


и~правоустанавливающих документов в~об\-ласти AI.


\end{itemize}


     


     \textbf{Второй этап (до 2025~г.)}: достижение прорыва в~базовой 


теории AI и~мирового лидерства в~этой области. Выведение AI в~качестве 


основной движущей силы индустриального и~экономического обновления 


Китая:


     \begin{itemize}


\item достижение прорыва в~базовой тео\-рии и~технологии AI (воз\-мож\-ность 


независимого обуче\-ния);


\item широкое применение AI в~интеллектуальном производстве, 


интеллектуальных здравоохранении, градостроительстве, сельском 


хозяйстве, обороне и~др.; расширение финансовой базы индустрии AI 


с~400~млрд до~5~трлн юаней;


\item дальнейшее совершенствование законодательной базы в~об\-ласти~AI.


     \end{itemize}


     


     \textbf{Третий этап (до 2030~г.)}: выведение тео\-рии, технологии 


и~применения AI в~КНР на такой уровень, чтобы стать основным мировым 


центром в~этой области, обеспечение на этой основе лидерства Китая:


     \begin{itemize}


\item создание основы теории и~технологий AI нового поколения, 


достижение прорыва в~области человеческого интеллекта, независимого AI, 


смешанного и~группового интеллекта;


\item достижение лидерства в~мировой конкурентной индустрии AI, широкое 


внед\-ре\-ние AI в~производство и~жизнь, социальную об\-ласть, оборону; 


расширение финансовой базы индустрии AI с~1~трлн до~10~трлн юаней;


\item создание лидирующей тренировочной базы об\-слу\-жи\-ва\-юще\-го персонала 


AI, разработка корпуса законов в~об\-ласти~AI.


\end{itemize}


     


     В~Плане декларированы ключевые задачи, решение которых обеспечит 


его реализацию, показанные следующим классификатором.


\begin{enumerate}[1.]


\item \textbf{Создание общей открытой технологической системы AI, 


предусматривающей кооперацию.}


\begin{enumerate}[{1}.1]


\item \textit{Разработка нового поколения базовой теории систем AI}.


\begin{enumerate}[{1.1}.1]


\item Интеллектуальная теория обработки больших данных.


\item Вычислительная теория восприятия разнородных данных.


\item Развитая теория гибридного интеллекта.


\item Теория группового интеллекта.


\item Автономное совместное управление, оптимизация и~принятие решений.


\item Теория развитого машинного обучения.


\item Теория интеллектуальных вычислений мозга.


\item Теория квантовых интеллектуальных вычислений.


\end{enumerate}


\item \textit{Создание нового поколения общих ключевых технологических сис\-тем~AI.}


\begin{enumerate}[{1.2}.1]


\item Вычислительные <<движки>> (computing engines) знаний и~сервисные технологии 


знаний.


\item Технологии логических заключений на основе анализа разнородной информации.


\item Ключевые технологии группового интеллекта.


\item Улучшенные технологии и~архитектура гибридного интеллекта.


\item Человеконезависимые и~соответствующие интеллектуальные технологии.


\item Технологии виртуальной реальности и~моделирования.


\item Разработка элементной базы для~AI.


\item Обработка естественных языков.


\end{enumerate}


\item \textit{Макет инновационной платформы~AI.}


\begin{enumerate}[{1.3}.1]


\item Аппаратно-программная платформа.


\item Платформа сервисов группового интеллекта.


\item Платформа поддержки улучшенного гибридного интеллекта.


\item Платформа поддержки человеконезависимых сис\-тем.


\item Платформа обеспечения безопас\-ности.


\end{enumerate}


\item \textit{Ускорение подготовки кадров для разработки~AI.}


\begin{enumerate}[{1.4}.1]


\item Создание высококвалифицированной команды (со\-во\-куп\-ность коллективов 


рассматривается как единое целое).


\item Подготовка образованных и~талантливых кад\-ров.


\item Модернизация системы образования.


\end{enumerate}


\end{enumerate}


\item \textbf{Развитие эффективной интеллектуальной экономики.}


\begin{enumerate}[{2}.1]


\item \textit{Ускорение промышленного применения~AI.}


\begin{enumerate}[{2.1}.1]


\item Внедрение интеллектуального аппаратного и~программного обеспечения.


\item Интеллектуальные роботы.


\item Интеллектуальные средства передвижения.


\item Средства виртуальной и~дополненной ре\-аль\-ности.


\item <<Умные>> терминалы.


\item Фундаментальные компоненты интернета вещей.


\end{enumerate}


\item \textit{Содействие интеллектуальной индустриальной модернизации.}


\begin{enumerate}[{2.2}.1]


\item Интеллектуальное производство.


\item Интеллектуальное сельское хозяйство.


\item Интеллектуальная логистика.


\item Умные финансы.


\item Использование AI в~бизнес-ана\-ли\-ти\-ке.


\item Умный дом.


\end{enumerate}


\item \textit{Развитие концепции интеллектуального предприятия.}


\begin{enumerate}[{2.3}.1]


\item Широкомасштабная реклама и~содействие преобразованию в~интеллектуальное 


предприятие.


\item Содействие внедрению умных фабрик.


\item Культивирование соревновательности за звание AI-ли\-де\-ров в~различных областях 


индустрии.


\end{enumerate}


\item \textit{Обеспечение режима благоприятствования при внедрении~AI.}


\begin{enumerate}[{2.4}.1]


\item Использование региональных особенностей и~преимуществ.


\item Создание пилотных участков.


\item Создание национального парка~AI.


\item Создание национальной базы~AI.


\end{enumerate}


\end{enumerate}


\item \textbf{Создание безопасного и~комфортного интеллектуального общества.}


\begin{enumerate}[{3}.1]


\item \textit{Разработка комфортных и~эффективных интеллектуальных сервисов.}


\begin{enumerate}[{3.1}.1]


\item Образование.


\item Здравоохранение.


\item Забота о престарелых.


\end{enumerate}


\item \textit{Продвижение AI в~социальном управ\-ле\-нии.}


\begin{enumerate}[{3.2}.1]


\item Интеллектуальное правительство.


\item Справедливое и~мудрое судопроизводство.


\item Умные города.


\item Интеллектуальная транспортная сис\-тема.


\end{enumerate}


\item\textit{Внедрение AI для укрепления общественной безопас\-ности.}


\item \textit{Содействие социальному общению и~доверию.}


\end{enumerate}


\item \textbf{Укрепление оборонной мощи.}


\item \textbf{Построение безопасной и~эффективной интеллектуальной 


инфраструктуры.}


\begin{enumerate}[{5}.1]


\item \textit{Коммуникационная инфраструктура}.


\item \textit{Big Data~--- инфраструктура}.


\item \textit{Высокопроизводительная вычислительная инфраструктура}.


\end{enumerate}


\item \textbf{Совместная разработка проектов в~других областях науки и~AI.}


\end{enumerate}





     В~Плане декларированы принципы распределения ресурсов при его 


реализации:


     \begin{itemize}


\item обеспечение финансовых приоритетов, использование рыночных 


механизмов;


\item оптимизация финансовой структуры при разработках AI;


\item координация международных и~национальных инновационных ресурсов;


\item обеспечение мер доверия.


\end{itemize}


     


     Помимо официального Плана~\cite{6-b} интерес вызывают 


иностранные (прежде всего американские) статьи, посвященные китайскому 


стилю реализации государственной политики в~области AI R\&D, например 


работа~\cite{8-b}. В~\cite{8-b} выражается обеспокоенность по поводу 


неограниченного использования AI в~китайских вооруженных силах, что дает 


КНР серьезное преимущество. Тем не менее это не мешает частным 


американским компаниям сотрудничать с~китайскими исследователями 


в~разработке технологий на основе AI.


     


     Одной из активно сотрудничающих с~китайскими исследователями 


компаний является Microsoft, которая работает с~фи\-нан\-си\-ру\-емым 


вооруженными силами Национальным университетом оборонных 


технологий (NUDT) над технологией наблюдения. Другой гигант, Google, 


открыл Исследовательский центр по искусственному интеллекту в~Пекине 


в~2017~г.\ Лаборатория разрабатывает инструменты для AI и~машинного 


обучения, вклю\-чая платформу с~открытым исходным кодом TensorFlow. Эти 


компании не несут от\-вет\-ст\-вен\-ности за воз\-мож\-ное использование их 


технологий в~будущем.


     


     Еще одним примером служит компания DJI (Шеньчжэнь, КНР), 


специализирующаяся на производстве дронов и~мультикоптеров. Дело в~том, 


что программное обеспечение для управ\-ле\-ния этими устройствами 


разрабатывается в~дивизионе DJI, расположенном в~Пало-Аль\-то  


(Сан\-та-Кла\-ра, Калифорния, США) американскими специалистами.


     


     Еще одним примером технологического сотрудничества служат 


совместные работы фирмы Baidu (КНР) и~Intel (Сан\-та-Кла\-ра, Калифорния, 


США) по созданию высокопроизводительных специализированных 


 AI-чи\-пов.


     


     Согласно исследованиям университета Синьхуа, примерно половина 


высокоцитируемых научных статей в~области AI R\&D написана китайскими 


исследователями совместно с~учеными из других стран.


     


\section{Структура финансирования AI R\&D в~КНР}


     


     Финансирование НИОКР в~сфере AI R\&D в~КНР обеспечивается 


тремя источниками~\cite{9-b}: центральным правительством, 


территориальными органами государственной власти  


и~фир\-ма\-ми\!/\!пред\-при\-яти\-ями.


     


     Структура правительства КНР, Госсовета, во многом схожа со 


структурой правительств РФ и~США. В~состав Госсовета входят такие 


профильные подразделения, как Министерство науки и~технологии 


и~Министерство промышленности и~информатизации. Распределение 


центрального государственного финансирования осуществляется через 


данные профильные министерства и~подчиненные им подразделения 


и~органы.


     


     К~территориальным органам власти КНР относятся правительства 


провинций, крупных городов, муниципальных образований и~пр. Обычно их 


финансирование AI R\&D не является значимым: ка\-кие-ли\-бо относительно 


специфические программы финансируются некоторыми территориальными 


органами власти.


     


     Частные предприятия в~Китае больше похожи на аналогичные 


экономические структуры в~РФ, чем в~США. Дело в~том, что в~КНР имеется 


значительная часть предприятий со смешанной  


част\-но-го\-су\-дар\-ст\-вен\-ной формой собственности (ЧГС). Хотя 


предприятия ЧГС, как правило, стремятся к~получению прибыли, участие 


государства создает стимулы для поддержки целей государственной 


политики, которые преобладают над целью получения прибыли. Кроме того, 


даже предприятия с~полностью част\-ной формой соб\-ст\-вен\-ности подвержены 


значительному влиянию со стороны правительства Китая, и~уровень такого 


влияния значительно выше, чем в~западных странах. Тот факт, что 


государственные деньги и~влияние тес\-но связаны с~деятельностью многих 


китайских предприятий, позволяет понять, почему иногда расходы 


предприятий ЧГС могут в~ка\-кой-то мере рас\-смат\-ри\-вать\-ся как 


государственные расходы.


     


     В~2014~г.\ Госсовет КНР анонсировал реорганизацию инновационной 


системы КНР. Реформа предполагала к~концу 2018~г.\ консолидацию сотен 


пересекающихся, избыточных и~неэффективных на\-уч\-но-тех\-ни\-че\-ских 


инновационных программ в~5~основных публично открытых источников 


финансирования:


     \begin{enumerate}[(1)]


\item Национальный естественно-научный фонд Китая (НЕНФК), 


специализирующийся на фундаментальных и~прикладных исследованиях;


\item целевые национальные программы R\&D, ориентированные на решение 


фундаментальных проблем улучшения социального и~экономического 


бла\-го\-со\-сто\-яния населения КНР;


\item национальные научно-технологические мегапроекты, направленные на 


достижение прорывов в~ключевых технологиях;


\item управляющие фонды, предостав\-ля\-ющие средства стартапам 


и~небольшим предприятиям для стимулирования инноваций и~передачи 


технологий;


    \item программа <<Основа и~таланты>> для поощрения создания 


исследовательских цент\-ров и~привлечения талантливых ученых для работы 


в~КНР.


\end{enumerate}





     Следует отметить, что в~этом перечне нет источников финансирования 


AI R\&D в~области обороны и~безопасности. При этом анализ направлений 


исследований позволяет заключить, что со\-от\-вет\-ст\-ву\-ющие результаты AI 


R\&D допускают двойное (гражданское и~военное) применение.


     


     НЕНФК финансирует исследования по следующим направлениям, 


связанным с~AI:


     \begin{itemize}


\item теоретические основания AI;


\item машинное обучение;


\item машинное восприятие и~распознавание образов;


\item обработка естественных языков;


\item представление и~обработка знаний;


\item интеллектуальные сис\-те\-мы и~приложения;


\item когнитивный и~нейробиологический~AI;


\item информатика в~образовании;


\item математические и~информационные междисциплинарные задачи.


\end{itemize}


     


     Целевые национальные программы AI R\&D КНР имеют несколько 


иную направленность:


     \begin{itemize}


\item теоретические основания AI;


\item обработка и~анализ больших объемов данных;


\item сетевые физические системы с~поддержкой вычислений;


\item разработка крупномасштабных сетей;


\item интеллектуальные роботизированные и~автономные системы;


\item информационные системы в~здравоохранении;


\item продуктивность, устой\-чи\-вость и~качество программного обеспечения;


\item обеспечение НИОКР высокопроизводительными вы\-чис\-ли\-тель\-ны\-ми 


сис\-те\-мами;


\item создание человеко-машинных систем.


\end{itemize}


     


     Национальные научно-технологические мегапроекты обычно имеют 


сле\-ду\-ющую на\-прав\-лен\-ность:


     \begin{itemize}


\item теоретические основания AI;


\item создание человеко-ма\-шин\-ных сис\-тем;


\item интеллектуальные роботизированные и~автономные сис\-те\-мы;


\item обеспечение НИОКР высокопроизводительными вычислительными  


сис\-те\-мами.


\end{itemize}


     


     Для примера широты и~детальности AI R\&D, по\-пав\-ших в~2018~г.\ 


в~мегапроект КНР <<Искусственный интеллект нового поколения>>, можно 


привести сформулированные в~его рам\-ках задачи:


     \begin{itemize}


\item новое поколение моделей нейронных сетей;


\item адаптивное восприятие для открытых окружений;


\item кросс-медийные при\-чин\-но-след\-ст\-вен\-ные выводы;


\item принятие решений в~играх с~неполной информацией;


\item механизм возникновения и~метод расчета группового интеллекта;


\item повышение интеллектуальных возможностей че\-ло\-ве\-ко-ма\-шин\-ных 


комплексов;


\item коллективные методы управления и~тео\-рии принятия решений 


в~слож\-ной производственной среде;


\item обобщенный механизм обучения и~расчета знаний предметной об\-ласти;


\item система кросс-медиа-ана\-ли\-за и~рассуждений;


\item технология активного восприятия сцены с~когнитивными целями;


\item исследование мотивации и~конвергенции группового интеллекта, 


ориентированного на групповое поведение;


\item исследования в~об\-ласти программного и~аппаратного обеспечения для 


взаимодействия человека и~машины;


\item высокоточные интеллектуальные измерения и~манипуляции 


в~беспилотных сис\-те\-мах;


\item точное интеллектуальное обучение работе автономных агентов;


\item новые датчики и~чипы;


\item ключевые стандарты проверки процессоров нейронных сетей.


\end{itemize}


     


\section{Организация AI R\&D в~области обороны и~безопасности КНР}


     


     Экономические успехи КНР имеют прямое отношение к~национальной 


безопас\-ности Китая~\cite{10-b}. Во-пер\-вых, это снижает возможности для 


других стран по дипломатическому и~экономическому давлению.  


Во-вто\-рых, это увеличивает технологический потенциал, доступный 


китайскому военному и~разведывательному сообществу. Практически все 


крупные технологические компании в~Китае активно сотрудничают 


с~китайскими военными и~государственными службами безопасности 


и~обязаны это делать по закону. Статья~7 <<Закона о~национальной разведке 


КНР>> дает правительству полномочия для принуждения к~такой помощи, 


хотя у~правительства также есть действенные ненасильственные 


инструменты стимулирования сотрудничества в~этой об\-ласти.  


Во\-ен\-но-граж\-дан\-ская интеграция является одним из краеугольных 


камней национальной стратегии КНР в~об\-ласти AI R\&D. При этом следует 


иметь в~виду, что огромный объем исследований в~области AI R\&D 


с~результатами двойного назначения проводится как китайскими 


государственными организациями, так и~част\-ны\-ми компаниями.


     


     Под эгидой NUDT в~начале 2018~г.\ были созданы две круп\-ные 


исследовательские организации, \mbox{целью} которых ставится проведение AI 


R\&D в~об\-ласти обороны. Это Исследовательский центр беспилотных 


сис\-тем (USRC) и~Исследовательский центр искусственного интеллекта 


(AIRC). Данные организации не относятся к~крупным по мировым 


масштабам~[10]: их общий штат со\-став\-ля\-ет около~200~человек по сравнению 


с~600~сотрудниками компании SenseTime (Гонконг, КНР) или 


700~сотрудниками DeepMind (дивизион Google, США). Тем не менее следует 


подчеркнуть, что USRC и~AIRC являются организациями, 


спе\-ци\-а\-ли\-зи\-ру\-ющи\-ми\-ся исключительно на оборонных разработках  


в~об\-ласти AI.


     


     Среди направлений AI R\&D в~об\-ласти обороны КНР наиболее 


важ\-ны\-ми пред\-став\-ля\-ют\-ся~[11--17]:


    \begin{description} 


\item[\,] \textbf{создание автономных платформ.} Китайские ученые 


предполагают применение AI в~новых <<автономных плат\-фор\-мах военного 


назначения>>. Они включают в~себя спутники, дроны, беспилотные 


автомобили, наземные роботы, надводные и~подводные суда, 


интеллектуальные боеприпасы и~др.;


     \item[\,] \textbf{создание вспомогательных систем.} Применение AI 


в~специальных ап\-па\-рат\-но-про\-грам\-мных сис\-те\-мах, соз\-да\-ва\-емых 


в~интересах разведки, рекогносцировки, слежения и~глубокой обработки 


данных. Разработка программных сис\-тем, обеспечивающих кибератаки 


и~защиту от них, ап\-па\-рат\-но-про\-грам\-мных средств радиоэлектронной


борьбы (РЭБ). Разработка 


комплексов автоматического принятия решений при пусках ракет;


     


 \item[\,] \textbf{разработка человеко-ма\-шин\-ных комплексов военного 


назначения.} Раз\-ра\-бот\-ка обуча\-ющих комп\-лек\-сов\!/\!тре\-на\-же\-ров военного 


персонала различного уров\-ня, специализации и~принадлежности к~родам 


войск (ВС, ВВС, ВМФ, ПВО, \mbox{РВСН}, РЭБ и~пр.), средств планирования 


военных операций, средств обучения автономных платформ 


и~взаимодействия с~ними.


\end{description}


     


     В~сфере обеспечения общественной безопасности КНР является самой 


передовой державой в~об\-ласти раз\-ра\-бот\-ки\!/\!внед\-ре\-ния и~интеграции 


различных сис\-тем наблюдения (AI surveillance systems). Приведем некоторые 


наиболее развитые на\-прав\-ле\-ния~[18]:


 \begin{description}    


\item[\,] \textbf{технологии\!/\!системы распознавания лиц ({Facial 


Recognition Technologies\!/\!Systems}).} Это набор биометрических технологий, 


использующих фото- и~видеокамеры для со\-по\-став\-ле\-ния неподвижных 


изображений людей или видеотреков с~име\-ющи\-ми\-ся в~специальных базах 


данных. Не все сис\-те\-мы \mbox{распознавания} лиц направлены на индивидуальную 


идентификацию, некоторые предназначены для мониторинга и~оценки 


разнообразных демографических тенденций и~проведения анализа 


эмоционального со\-сто\-яния толпы с~по\-мощью распознавания отдельных 


лиц;


    \item[\,] \textbf{умный город/безопасный город (Smart City/Safe City).} 


Согласно~[19], это набор технологий и~приемов интеграции общества, 


экономики и~окру\-жа\-ющей среды, повышающий качество процессов во всех 


сферах человеческой жиз\-не\-дея\-тель\-ности, обес\-пе\-чи\-ва\-ющий устойчивое 


развитие, безопас\-ность и~здоровье жителей с~\mbox{целью} повышения качества 


жизни граждан. С~точки зрения безопас\-ности эта концепция предполагает 


включение в~себя сле\-ду\-ющих <<умных>> (т.\,е.\ использующих технологии 


AI) функциональных подсистем: транспортной подсистемы и~оптимизации 


маршрутов; системы наблюдения, поиска, обнаружения и~идентификации; 


системы управления в~кризисных ситуациях, их предсказания, раннего 


обнаружения, предупреждения и~мониторинга; централизованной сис\-те\-мы 


обеспечения оперативной деятельности органов правопорядка и~экстренного 


реагирования; внед\-ре\-ния технологий интернета вещей, обеспечения доступа к~интернету, 


обеспечения защиты данных и~пр.;


     \item[\,] \textbf{умное обеспечение правопорядка (Smart Policing, SP).} 


Идея SP заключается в~агрегировании и~на\-стра\-ива\-емой автоматической 


обработке большого объема разнородных данных: географических, 


исторических, криминальных, технических, био\-мет\-ри\-че\-ских  


и~социальных~--- для предотвращения преступлений, оперативной реакции 


на них и~их прогнозирования.


\end{description}


     


     Все эти технологии обеспечения безопасности широко внед\-ре\-ны 


     и~продолжают тестироваться и~развиваться по всей территории КНР. Помимо 


этого, различные со\-став\-ные час\-ти этих сис\-тем и~сис\-те\-мы в~целом 


по\-став\-ля\-ют\-ся IT-ли\-де\-ра\-ми КНР~--- фирмами Huawei, ZTE, Hikvision, 


SenseTime, Medvii, Yitu~--- более чем в~50~стран.





\vspace*{-6pt}





\section{Заключение}





\vspace*{-2pt}


     


     Завершая в~этой статье состоящее из четырех частей исследование 


влияния государственного управления на сферу AI R\&D, отметим, во-пер\-вых, 


что относительно государственной политики в~этой об\-ласти 


существует весьма обширный аналитический материал: достаточно 


упомянуть работы~[20--23]. И~в~итоге анализа этих и~других работ по 


тематике цикла приведем сле\-ду\-ющие пред\-став\-ля\-ющи\-еся наиболее важными 


выводы.


     \begin{enumerate}[1.]


\item Все государства, независимо от уровня их промышленного развития, 


характеризуют перспективы внедрения технологий AI как ключевой фактор, 


способный повлиять и~на уровень жизни в~различных странах, и~на 


изменение межгосударственной табели о рангах, мирового промышленного, 


политического и~военного лидерства.


\item Промышленно развитые страны, а~так\-же государства~--- региональные 


лидеры, об\-ла\-да\-ющие политическими амбициями и~достаточными  


на\-уч\-но-тех\-ни\-че\-ски\-ми и~материальными ресурсами, разработали 


стратегические документы, рег\-ла\-мен\-ти\-ру\-ющие пути развития AI R\&D 


и~внед\-ре\-ния соответствующих технологий в~своих странах, а также степень 


учас\-тия государственных органов в~этих процессах.


\item Стратегии развития AI R\&D в~разных странах можно 


классифицировать по степени полноты охвата на\-прав\-ле\-ний и~не\-за\-ви\-си\-мости 


проводимых исследований следующим образом.


\begin{enumerate}[3.1]


\item Стратегии, предполагающие полный спектр работ в~области AI R\&D 


и~декларирующие в~качестве одной из целей достижение лидерских позиций 


данного государства в~этой об\-ласти в~мире. Примерами таких государств 


могут служить США, КНР, РФ.


\item  Стратегии, предполагающие получение новых на\-уч\-но-тех\-ни\-че\-ских 


результатов в~процессе тесного межгосударственного сотрудничества и~на\-уч\-но-тех\-ни\-че\-ской 


кооперации. Примерами подобных государств могут 


служить страны ЕС.


\item Стратегии <<ключевых>> союзников. Они подразумевают 


сотрудничество с~США в~области AI R\&D, а~также признание их 


лидирующего положения. В~то же время <<ключевые>> союзники обладают 


широкой самостоятельностью по выбору на\-прав\-ле\-ний и~методов AI R\&D, 


а~также воз\-мож\-ности привлечения к~партнерству других государств. 


Примерами подобного <<ключевого>> союзника являются Израиль, а~также 


Великобритания, тес\-но со\-труд\-ни\-ча\-ющая с~США и~опи\-ра\-юща\-яся в~AI R\&D 


не только на свои научные ресурсы, но и~на ресурсы всего Британского 


Содружества.


\item Стратегии, более сфокусированные на внед\-ре\-ние уже су\-ще\-ст\-ву\-ющих 


<<чужих>> AI технологий, с~малым вниманием к~собственным R\&D.


\end{enumerate}


\end{enumerate}





\vspace*{-6pt}





{\small\frenchspacing


{%\baselineskip=10.8pt


\addcontentsline{toc}{section}{Литература}


\begin{thebibliography}{99}





\vspace*{-2pt}





\bibitem{1-b}


\Au{Борисов А.\,В., Босов~А.\,В., Жуков~Д.\,В.} Стратегия исследований и~разработок 


в~области искусственного интеллекта I: Основные понятия и~краткая хронология~/\!\!/ 


Сис\-те\-мы и~средства информатики, 2021. Т.~31. №\,1. С.~57--68.


\bibitem{2-b}


\Au{Борисов А.\,В., Босов~А.\,В., Жуков~Д.\,В.} Стратегия исследований и~разработок 


в~об\-ласти искусственного интеллекта II: Сравнительный анализ наукометрических 


показателей в~мире и~в Российской Федерации~/\!\!/ Сис\-те\-мы и~средства информатики, 2021. 


Т.~31. №\,2. С.~89--107.


\bibitem{3-b}


\Au{Борисов А.\,В., Босов~А.\,В., Жуков~Д.\,В.} Стратегия исследований и~разработок 


в~об\-ласти искусственного интеллекта~III: Доктрина государственной поддержки США~/\!\!/ 


Сис\-те\-мы и~средства информатики, 2021. Т.~31. №\,4. С.~114--134.


\bibitem{4-b}


Национальная стратегия развития искусственного интеллекта на период до 


2030~года~/\!\!/ Указ Президента РФ от 10~октября 2019~г. №\,490 <<О~развитии 


искусственного интеллекта в~Российской Федерации>>.


\bibitem{5-b}


The National Artificial Intelligence Research and Development Strategic Plan.~---  National 


Science and Technology Council, Networking and Information Research and Development 


Sub-committee, 2016. 48~p.


\bibitem{6-b}


Notice of the State Council Issuing the New Generation of Artificial Intelligence 


Development Plan: State Council Document [2017] No.\,35~/\!\!/ The Foundation for Law and 


International Affairs, 2017. 28~p. 


{\sf https://flia.org/wp-content/uploads/2017/07/A-New-Generation-of-Artificial-Intelligence-Development-Plan-1.pdf}.


\bibitem{7-b}


New Generation of Artificial Intelligence Development Plan~/\!\!/ China Science  Technology 


Newsletters, July~20, 2017. {\sf https://www.mfa.gov.cn/ce/cefi/eng/kxjs/\linebreak P020171025789108009001.pdf}.


\bibitem{8-b}


\Au{Abadicio M.} Artificial Intelligence in the Chinese military~--- current initiatives~/\!\!/ Emerj 


Artificial Intelligence Research, 2019.


{\sf  https://www.orfonline.org/research/\linebreak a-i-in-the-chinese-military-current-initiatives-and-the-implications-for-india-61253/}.


\bibitem{9-b}


\Au{Colvin~T.\,J., Liu~I., Babou~T.\,F., Wong~G.\,J.}


 A~brief examination of Chinese government expenditures on Artificial 


Intelligence R\&D.~--- FFRDC: Science and Technology Policy 


Institute, 2020.  IDA document D-12068. {\sf https://www.jstor.org/ stable/resrep22826}. 


\bibitem{10-b}


\Au{Allen G.\,C.} Understanding China's AI strategy: Clues to Chinese strategic thinking on 


Artificial Intelligence and national security.~--- Center for a~New American Security, 2019.


 {\sf https://www.cnas.org/publications/reports/understanding-chinas-ai-strategy}.


 


 \bibitem{13-b} %11


\Au{Kania E.\,B.} Chinese military innovation in Artificial Intelligence: Testimony before the 


U.S.--China economic and security review commission hearing on trade, technology, and 


military-civil fusion.~--- Center for a~New American Security, 2019. 53~p.


{\sf https://www.cnas.org/publications/congressional-testimony/chinese-military-innovation-in-artificial-intelligence}.


\bibitem{14-b} %12


\Au{Dahm M.} Chinese debates on the military utility of Artificial Intelligence~/\!\!/ War on the 


Rocks, 2020. {\sf https://warontherocks.com/2020/06/chinese-debates-on-the-military-utility-of-artificial-intelligence/}.





\bibitem{12-b} %13


\Au{Fedasiuk R.} Chinese perspectives on AI and future military capabilities.~--- Center for 


Security and Emerging Technology, 2020.


{\sf https://cset.georgetown.edu/wp-content/uploads/CSET-Chinese-Perspectives.pdf}.





\bibitem{11-b} %14


\Au{Tadjdeh Y.} China threatens U.S.\ primacy in Artificial Intelligence (updated)~/\!\!/ 


National Defence, 2020.


{\sf https://www.nationaldefensemagazine.org/articles/2020/ 10/30/china-threatens-us-primacy-in-artificial-intelligence}.





\bibitem{15-b}


\Au{Zhang J.} China's military employment of Artificial Intelligence and its security 


implications~/\!\!/ International Affairs Review, 2020. Vol.~XXVIII. Iss.~2. P.~38--56.


\bibitem{16-b}


\Au{Kania E.\,B.} ``AI weapons'' in China's military innovation~/\!\!/ Global China, 2020. 


23~p. {\sf https://www.brookings.edu/wp-content/uploads/2020/04/FP\_20200427\_ai\_


weapons\_kania\_v2.pdf}.


\bibitem{17-b}


\Au{Roberts H., Cowls~J., Morley~J., Taddeo~M., Wang~V., Floridi~L.} 


The Chinese approach to Artificial Intelligence: An analysis of policy, 


ethics, and regulation~/\!\!/ AI \& SOCIETY, 2021. Vol.~36. P.~59--77.


\bibitem{18-b}


\Au{Feldstein S.} The global expansion of AI surveillance.~--- Washington, DC, USA: 


Carnegie Endowment for International Peace, 2019. 42~p.


\bibitem{19-b}


\Au{Lacin$\acute{\mbox{a}}$k M., Ristvej~J.} Smart city, safety and security~/\!\!/ Procedia 


Engineer., 2017. Vol.~192. P.~522--527.


\bibitem{20-b}


\Au{Groth O.\,J., Nitzberg~M., Zehr~D.} Comparison of national strategies to promote Artificial 


Intelligence. Part~1.~--- Berlin: Konrad-Adenauer-Stiftung e.V., 2019. 70~p.


\bibitem{21-b}


\Au{Groth O.\,J., Nitzberg~M., Zehr~D.} Comparison of national strategies to promote Artificial 


Intelligence. Part 2.~--- Berlin: Konrad-Adenauer-Stiftung e.V., 2019. 84~p.


\bibitem{22-b}


\Au{Van Roy~V.} AI watch~--- national strategies on Artificial Intelligence: A~European 


perspective in 2019.~--- Luxembourg: Publications Office of the European 


Union, 2020.  EUR 30102 EN. 90~p. doi: 10.2760\!/\!602843.


{\sf https://publications.jrc.ec.europa.eu/\linebreak 


repository/bitstream/JRC119974/national\_strategies\_on\_artificial\_intelligence\_final\_1.\linebreak pdf}.


\bibitem{23-b}


Artificial Intelligence in the Asia-Pacific region: Examining policies and strategies to 


maximise AI readiness and adoption.~--- London: International Institute of Communications, 


2020. 88~p. {\sf https://www.iicom.org/wp-content/uploads/IIC-AI-Report-2020.pdf}.


       \end{thebibliography}


} }











\vspace*{-10pt}





\hfill{\small\textit{Поступила в~редакцию 21.04.21}}





\vspace*{6pt}





\hrule





\vspace*{2pt}





%\newpage





%\vspace*{-28pt}





\hrule





\vspace*{-2pt}





\def\tit{RESEARCH AND DEVELOPMENT STRATEGY\\ IN~THE~FIELD 


OF~ARTIFICIAL INTELLIGENCE~IV:\\ CHINESE GOVERNMENT POLICY\\[-5pt]}








\def\titkol{Research and development strategy in~the~field 


of~artificial intelligence~IV: Chinese %Government 


policy}





\def\aut{A.\,V.~Borisov, A.\,V.~Bosov, and D.\,V.~Zhukov}





\def\autkol{A.\,V.~Borisov, A.\,V.~Bosov, and D.\,V.~Zhukov}





\titel{\tit}{\aut}{\autkol}{\titkol}





\vspace*{-18pt}





 \noindent


     Federal Research Center ``Computer Science and Control'' of the Russian 


Academy of Sciences, 44-2~Vavilov Str., Moscow 119333, Russian Federation








\def\leftfootline{\small{\textbf{\thepage}


\hfill Sistemy i Sredstva Informatiki~---


Systems and Means of Informatics 2022 vol~32 no\,1}


}%


 \def\rightfootline{\small{Sistemy i Sredstva Informatiki~---


 Systems and Means of Informatics 2022 vol~32 no\,1


\hfill \textbf{\thepage}}}





\vspace*{1pt}





     


     


     


     


     \Abste{The paper concludes a cycle of works 


     devoted to the analysis of the impact of state-shock management on the effectiveness 


     of research and development in the field of artificial intelligence (AI R\&D). 


     The fourth part of the cycle presents the result of the analysis of the impact of the 


     state policy of the People's Republic of China (PRC) on the AI R\&D area. 


     Brief information on the scientific and technical potential and strategic 


     documents of the PRC in terms of the development of AI R\&D, adopted since 2017,


      is presented. The structure of public and private financing of AI R\&D in the PRC 


      is described. A~brief description of the organization of AI R\&D in the field of 


      defense and security is given. Concluding remarks are presented including 


      differences in approaches to strategic


      state planning in the field of AI R\&D adopted in different states.}


     


     \KWE{artificial intelligence; face recognition system; 


surveillance/tracking system; smart city/safe city}





     


\DOI{10.14357\!/\!08696527220102} 





%\vspace*{-16pt}








 %\Ack


 % \noindent


 








%\vspace*{-6pt}





\renewcommand{\bibname}{\protect\rmfamily References}


%\renewcommand{\bibname}{\large\protect\rm References}





\vspace*{-16pt}





{\small\frenchspacing


{%\baselineskip=10.8pt


\addcontentsline{toc}{section}{References}


\begin{thebibliography}{99}





\vspace*{-4pt}





     \bibitem{1-b-1}


     \Aue{Borisov, A.\,V., A.\,V.~Bosov, and D.\,V.~Zhukov.} 2021. Strategiya 


issledovaniy i~raz\-ra\-bo\-tok v~oblasti iskusstvennogo intellekta I: Osnovnye 


ponyatiya i~kratkaya khro\-no\-lo\-giya [Reseach and development strategy in the field 


of artificial intelligence I: Basic concepts and brief chronology]. \textit{Sistemy 


i~Sredstva Informatiki~--- Systems and Means of Informatics} 31(1):57--68.


 \bibitem{2-b-1}


     \Aue{Borisov, A.\,V., A.\,V.~Bosov, and D.\,V.~Zhukov.} 2021. Strategiya 


issledovaniy i~raz\-ra\-botok v~oblasti iskusstvennogo intellekta II: Sravnitel'nyy 


analiz na\-uko\-met\-ri\-che\-skikh pokazateley v~mire i~v~Rossiyskoy Federatsii [Reseach 


and development strategy in the field of artificial intelligence II: Comparative 


analysis of scientometric\linebreak\vspace*{-12pt}





\pagebreak





\noindent


 indicators in the world and in the Russian Federation]. 


\textit{Sistemy i~Sredstva Informatiki~--- Systems and Means of Informatics} 31(2):89--107.


\bibitem{3-b-1}


     \Aue{Borisov, A.\,V.,  A.\,V.~Bosov, and D.\,V.~Zhukov.} 2021. Strategiya 


issledovaniy i~raz\-ra\-bo\-tok v~oblasti iskusstvennogo intellekta III: Doktrina 


gosudarstvennoy pod\-derzh\-ki SShA [Research and development strategy in the 


field of artificial intelligence III: United States government support doctrine].


 \textit{Sistemy i~Sredstva Informatiki~--- Systems and Means of Informatics} 31(4):114--134.


\bibitem{4-b-1}


O razvitii iskusstvennogo intellekta v~Rossiyskoy Federatsii: ukaz 


Prezidenta ot 10.10.2019 No.\,490 [About strategy of scientific and technological 


development of the Russian Federation. Presidential Decree No.\,490 dated 


10.10.2019]. Available at: 


{\sf http://static.kremlin.ru/media/events/$\vert$les/ru/AH4x6HgKWANwVtMOfPDhcbRpvd\linebreak 1HCCsv.pdf}


 (accessed February~2, 2022).


\bibitem{5-b-1}


     National Science and Technology Council, Networking and 


Information Research and Development Sub-committee. 2016. The National 


Artificial Intelligence Research and Development Strategic Plan. Available at: 


{\sf www.nitrd.gov/pubs/\linebreak national\_ai\_rd\_strategic\_plan.pdf} (accessed February~2, 


2022).


\bibitem{6-b-1}


 Notice of the State Council Issuing the New Generation of Artificial Intelligence Development Plan. 2017. 


The Foundation for Law and International Affairs. 28~p. Available at: 


{\sf https://flia.org/wp-content/uploads/2017/07/A-New-Generation-\linebreak of-Artificial-Intelligence-Development-Plan-1.pdf}


 (accessed February~2, 2022).


\bibitem{7-b-1}


     New Generation of Artificial Intelligence Development Plan. July~20, 2017.


     \textit{China Science  Technology 


Newsletters}. Available 


at: {\sf https://www.mfa.gov.cn/ce/cefi/\linebreak eng/kxjs/P020171025789108009001.pdf} 


(accessed February~2, 2022).


\bibitem{8-b-1}


\Aue{Abadicio, M.} 2019.


     Artificial Intelligence in the Chinese military~--- current initiatives. 


     \textit{Emerj Artificial Intelligence Research}.


Available at: 


{\sf  https://www.orfonline.org/\linebreak research/a-i-in-the-chinese-military-current-initiatives-and-the-implications-for-india-61253/}


 (accessed February~2,  2022). 


 





   \bibitem{9-b-1}


   \Aue{Colvin, T.\,J., I.~Liu, T.\,F.~Babou, and G.\,J.~Wong.} 2020.


A brief examination of Chinese government expenditures on Artificial 


Intelligence R\&D.  FFRDC: Science and Technology Policy 


Institute. IDA document D-12068. Available at: {\sf https://www.jstor.org/stable/resrep22826} 


(accessed February~2, 2022).


\bibitem{10-b-1}


\Aue{Allen, G.\,C.} 2019.


Understanding China's AI strategy: Clues to Chinese strategic thinking 


on Artificial Intelligence and national security.  Center for a~New American Security. Available at: 


{\sf https://www.cnas.org/publications/reports/understanding-chinas-ai-strategy}


(accessed February~2, 2022).





\bibitem{13-b-1} %11


\Aue{Kania, E.\,B.} 2019.


     Chinese military innovation in Artificial Intelligence: Testimony before 


the U.S.--China economic and security review commission hearing on trade, 


technology, and military-civil fusion. Center for a~New American Security. 53~p. Available at: 


{\sf https://www.cnas.org/publications/congressional-testimony/\linebreak chinese-military-innovation-in-artificial-intelligence}


 (accessed February~2, 2022).


\bibitem{14-b-1} %12


\Aue{Dahm, M.} 2020.


Chinese debates on the military utility of Artificial Intelligence.


\textit{War on the Rocks}.  


Available at: 


{\sf https://warontherocks.com/2020/06/chinese-debates-on-the-military-utility-of-artificial-intelligence/}


 (accessed February~2, 2022).


 


 \bibitem{12-b-1} %13


\Aue{Fedasiuk, R.} 2020.


Chinese perspectives on AI and future military capabilities. Center for Security and Emerging Technology. 


Available at: 


{\sf https://cset.georgetown.edu/wp-content/uploads/CSET-Chinese-Perspectives.pdf}


 (accessed February~2, 2022).


 


\bibitem{11-b-1} %14


    \Aue{Tadjdeh Y.} 2020. China threatens U.S.\ primacy in Artificial Intelligence (updated). 


    \textit{National Defence}.


 Available at: 


{\sf https://www.nationaldefensemagazine.org/\linebreak articles/2020/10/30/china-threatens-us-primacy-in-artificial-intelligence}


 (accessed February~2, 2022).








\bibitem{15-b-1}


\Aue{Zhang, J.} 2020. China's military employment of artificial intelligence 


and its security implications. \textit{International Affairs Review} XXVIII(2):38--56.


\bibitem{16-b-1}


\Aue{Kania, E.\,B.} 2020.


``AI weapons'' in China's military innovation. \textit{Global China}. 23~p. Available at: 


{\sf https://www.brookings.edu/wp-content/uploads/2020/04/\linebreak  FP\_20200427\_ai\_weapons\_kania\_v2.pdf} (accessed 


February~2, 2022).


\bibitem{17-b-1}


\Aue{Roberts, H., J.~Cowls, J.~Morley, M.~Taddeo, V.~Wang, and L.~Floridi.} 


2021. The Chinese approach to Artificial Intelligence: An analysis of policy, ethics, 


and regulation. \textit{AI \& SOCIETY} 36:59--77.


\bibitem{18-b-1}


\Aue{Feldstein, S.} 2019. \textit{The global expansion of AI surveillance}. Washington, 


DC: Carnegie Endowment for International Peace. 42~p.


\bibitem{19-b-1}


\Aue{Lacin$\acute{\mbox{a}}$k, M., and J.~Ristvej.} 2017. Smart city, safety 


and security. \textit{Procedia Engineer.} 192:522--527.


\bibitem{20-b-1}


\Aue{Groth, O.\,J., M.~Nitzberg, and D.~Zehr.} 2019. 


\textit{Comparison of national 


strategies to promote Artificial Intelligence. Part~1.} Berlin: Konrad-Adenauer-Stiftung e.V. 70~p. 


\bibitem{21-b-1}


\Aue{Groth, O.\,J., M.~Nitzberg, and D.~Zehr.} 2019. \textit{Comparison of national 


strategies to promote Artificial Intelligence. Part~2.} Berlin: Konrad-Adenauer-Stiftung e.V. 84~p.


\bibitem{22-b-1}


\Aue{Van Roy, V.} 2020.


AI watch~--- national strategies on Artificial Intelligence: A~European 


perspective in 2019. Luxembourg: Publications Office of the European Union. EUR 30102 EN.


90~p. doi: 10.2760\!/\!602843.


Available at: 


{\sf https://publications.jrc.ec.\linebreak


 europa.eu/repository/bitstream/JRC119974/national\_strategies\_on\_artificial\_\linebreak intelligence\_final\_1.pdf}


 (accessed February~2, 2022).


\bibitem{23-b-1}


Artificial intelligence in the Asia-Pacific region: Examining policies 


and strategies to maximise AI readiness and adoption. 2020. 


London: International Institute of Communications. 


88~p. Available at: 


{\sf https://www.iicom.org/wp-content/uploads/IIC-AI-Report-2020.pdf} (accessed 


February~2, 2022).


\end{thebibliography}


} }





\vspace*{-6pt}





\hfill{\small\textit{Received April 21, 2021}} 





\vspace*{-15pt}  


     


     \Contr


     


     \vspace*{-3pt}


     


      \noindent


     \textbf{Borisov Andrey V.} (b.\ 1965)~--- Doctor of Science in physics and 


mathematics, principal scientist, Institute of Informatics Problems, Federal 


Research Center ``Computer Science and Control'' of the Russian Academy of 


Sciences, 44-2~Vavilov Str., Moscow 119333, Russian Federation; 


\mbox{ABorosov@ipiran.ru}





     \vspace*{3pt}


     


     \noindent


     \textbf{Bosov Alexey V.} (b.\ 1969)~--- Doctor of Science in technology, 


principal scientist, Institute of Informatics Problems, Federal Research Center 


``Computer Science and Control'' of the Russian Academy of Sciences,  


44-2~Vavilov Str., Moscow 119333, Russian Federation; 


\mbox{AVBosov@ipiran.ru}


     


     


     \vspace*{3pt}


     


     \noindent


     \textbf{Zhukov Denis V.} (b.\ 1979)~--- principal specialist, Institute of 


Informatics Problems, Federal Research Center ``Computer Science and Control'' 


of the Russian Academy of Sciences, 44-2~Vavilov Str., Moscow 119333, Russian 


Federation; \mbox{DZhukov@ipiran.ru}





\label{end\stat}





\renewcommand{\bibname}{\protect\rm Литература} 


     


     


     


